% ===================================================================
% Section 11. Contexts, idealization, and checking physics
% Sources: research/theory/T3-contexts-idealization-export.md (as repaired; see its
%          "Verification log" and research/verification/reverification-round2.md),
%          T2 (Prop 7.1), T6 (Sec. 4), research/lit/L7, L9, L10.
% Proofs and secondary results: sections/app-physics.tex.
% ===================================================================
\providecommand{\fM}{\mathfrak{M}}
\providecommand{\fF}{\mathcal{F}}
\providecommand{\fG}{\mathcal{G}}
\providecommand{\fK}{\mathcal{K}}
\providecommand{\tset}[1]{\lVert #1\rVert}
\providecommand{\cder}{\Vdash}
\providecommand{\dom}{\operatorname{dom}}
\providecommand{\osc}{\operatorname{osc}}
\providecommand{\AUX}{\mathsf{AUX}}
\providecommand{\Anc}{\mathrm{Anc}}
\providecommand{\Ccert}{C_{\mathrm{cert}}}
\providecommand{\Asm}{\mathrm{Asm}}
\providecommand{\Ax}{\mathrm{Ax}}
\providecommand{\ACCEPT}{\textsf{ACCEPT}}
\providecommand{\REJECT}{\textsf{REJECT}}
\providecommand{\rl}[1]{\textup{(#1)}}

\section{Contexts, idealization, and checking physics}\label{sec:physics}

Physics is where the programme of the motivating question meets its hardest test, and
where the problem of contexts cannot be postponed. Existing olympiad solutions, the
motivating notes observe, often do not reduce the problem to one clear mathematical
problem: ``there's more of a dance with many clear setups, with some relations to the
physical problem of interest'' \citep[\texttt{ai/verification and generation/verification
of physics olympiad solutions/the structure of physics olympiad solutions.md}]{hanni2026notes}.
Some setups are suppositions made for reductio. Some are idealizations that contradict
known facts (``assume the air pressure is $0$''). Some complete a problem that leaves
something open (the sun's elevation is not given). A coherence loss that penalizes good
arguments for both $P$ and $\neg P$ would penalize all three. A checker that asks whether
each claim is true fails at once, because good solutions are full of literally false
claims; the motivating notes ask what else it could check: ``check that a claim is true in
the context at hand. but wtf is that???'' (\texttt{ai/verification and
generation/verification from truth.md}).

This section answers with a semantics of contexts (\S\ref{sec:physics:semantics}), a
theory of coherence, export and tolerances
(\S\S\ref{sec:physics:coherence}--\ref{sec:physics:learning}), and a checker that is sound
against every solver relative to a provably ineliminable judgment layer, with a worked run
(\S\S\ref{sec:physics:checker}--\ref{sec:physics:example}).

We build a \emph{checker}, not a solver. Gold-level scores on IPhO~2025 theory under
rubric grading have been reported for current AI systems
\citep{yu2025hipho,qiu2025supernova}, though in the benchmark's September 2025 data not
on the modelling-heavy EuPhO~2025 of our worked example (best model $14.9$ of $29$, gold
threshold $16.5$; a later update reports gold on all its olympiads). The unmet target is
principled checking. The motivating notes also hold that publishing insightful research
on AI theorem provers is gravely dangerous (``in my opinion, publishing such research is
omnicidal'', \texttt{ai/verification and generation/formalizing philosophy.pdf}, a July
2024 draft). We have therefore avoided solver-side contributions, and leave the decision
about publication to the author of the notes.

Most results here are trivial once the definitions are in place, or import standard tools
(\L{}o\'s's theorem, Gronwall's inequality, split conformal prediction, information-based
complexity, Lipschitz covering bounds, interval arithmetic). New, as far as we know, are the frame-free
realizability criterion (\cref{thm:physics:realizability}(a)), the two-sided analysis of
adversarial stipulation (\cref{thm:physics:stipulation}), the table of one-sided signals
(\cref{tab:physics:onesided}), and a rigorous thin-leg certificate with a complete worked
checker run (\cref{sec:physics:example}). Proofs not given inline are in
Appendix~\ref{app:physics}.

% -------------------------------------------------------------------
\subsection{The problem: eternalism, contexts, and three import regimes}\label{sec:physics:problem}

The simplest coherence constraint is \emph{eternalist}: each proposition is true or
false independently of context, and a learner is penalized whenever it can argue for
both $P$ and $\neg P$. Physics reasoning violates this constantly and correctly. A
hydrostatics problem stipulates $p_{\rm atm}=0$, while the background derives
$p_{\rm atm}\approx1.0\times10^5$\,Pa. A reductio derives $\bot$ under a supposition.
One exam treats the same physics in incompatible models: at IPhO~2025, Theory Problem~2
excludes surface tension and lets a gas phase appear exactly at the equilibrium
threshold, while Theory Problem~3 makes surface tension the mechanism of bubble growth.
The motivating notes record the difficulty with two obvious moves. Conjoining every
sentence with the whole background framework ``leads to just assigning probability $0$
to everything''; and proofs by contradiction inside a framework known to be inadequate
``might \dots\ really involve subverting the background framework''
(\texttt{philosophy/epistemology/probabilities/assigning probabilities in a conception of
the world you know to be inadequate.md}).

\paragraph{Three import regimes.} ``Suppose'' covers three operations that differ in
what the new context imports from its parent:
\begin{center}\small
\begin{tabularx}{\textwidth}{@{}l>{\raggedright\arraybackslash}X>{\raggedright\arraybackslash}X>{\raggedright\arraybackslash}X>{\raggedright\arraybackslash}X@{}}
\toprule
Regime & Example & Imports & $\bot$ inside means & What leaves the context\\
\midrule
supposition & ``suppose $\sqrt2=p/q$'' & everything & success (reductio) & $\neg A$, or $A\to B$\\
counterfactual & ``if the atmosphere vanished, the water would boil'' & a maximal compatible part & the revision was done badly & $A\mathrel{\Box\!\!\to}B$\\
idealization & ``neglect drag'', ``$p_{\rm atm}=0$'', ``massless rope'' & a designated minimal fragment & the model is ill-posed & query answers, through a certified bridge\\
\bottomrule
\end{tabularx}
\end{center}
Physics idealizations are \emph{not} counterfactuals. ``If the air pressure were $0$,
the water in the beaker would boil'' is true and irrelevant: the problem asks what follows
in a designated model that imports incompressible hydrostatics and the data, and nothing
else. Similarity semantics
\citep{stalnaker1968theory,lewis1973counterfactuals} or AGM revision
\citep{alchourron1985logic} imports exactly the background (the vapour pressure of
water) that makes the context degenerate. Idealization is designated minimal import,
as in chunk-and-permeate \citep{brown2004chunk}, McCarthy's $\ist(c,\varphi)$
\citep{mccarthy1993notes}, local-models semantics \citep{ghidini2001local}, Cartwright's
models \citep{cartwright1983laws}, and compositional modelling
\citep{falkenhainer1991compositional}. We keep their shape, ``classical inside,
filtered between''; what none of them supplies, as far as we know, is different coherence
obligations for suppositions and idealizations, and rigorous certificates on the bridges.

\paragraph{Anatomy of a solution.} Three decomposed olympiad solutions (IPhO~2012
Theory Problem~1, part~A; EuPhO~2025 Theory Problem~1(a); IPhO~2025 Theory Problem~2)
share one structure. A \emph{root} $c_P$ holds the problem text, figures and
conventions; for grading, the ``world'' is $c_P$, not the actual world. A
\emph{reading} $\rho$ maps $c_P$ to an admissible model family, possibly with free
completion parameters, and types the query (``express in terms of $a,I_0,r,\varphi$'').
\emph{Local models} are clear mathematical problems, joined by \emph{bridges}: exact
ones (symmetry, decomposition, isomorphism) and approximate ones (a small parameter
set to its limit, with an error claim), ending in an \emph{export}. Only approximate
bridges introduced by the solver, rather than stipulated or regime-stated by the problem,
need a certificate from the solver. The notes' hypothesis that ``there's a local
setting up of a clear thing'' holds in all three: the unclear part concentrates in the
reading and in a few approximate bridges.

% -------------------------------------------------------------------
\subsection{Filter semantics for contexts}\label{sec:physics:semantics}

\begin{definition}[Frames, worlds, readings]\src{T3 Def 1.1}\label{def:physics:frame}
Let $L$ be a many-sorted first-order language with a sort $\R$ and parameter constants
$p_1,\dots,p_n$ of that sort. (i)~A \emph{deformation frame} is a partial map
$\fM:\Lambda\rightharpoonup\mathrm{Str}(L)$, $\Lambda\subseteq\R^n$, with
$p_i^{\fM(\lambda)}=\lambda_i$; $\dom\fM$ is where the model exists (is well-posed). In
every $\fM(\lambda)$ the sort $\R$ is the ordered field of reals with standard numerals;
$\fK$ is the class of $L$-structures with this standard real sort. Ideal values at
infinity are reparametrized ($m\to\infty$ as $1/m\to0$). (ii)~A \emph{world} is
$w=(\fM_w,\lambda^*_w)$ with $\lambda^*_w\in\dom\fM_w$: a top-model family and the actual
(or intended) parameter point. (iii)~A \emph{reading} $\rho$ yields a set $W_\rho$ of
worlds, the \emph{admissible completions} (given parameters fixed, measured ones in
intervals, free ones in ranges); a query term $Q$; and a set $\Ax_\rho$ of sentences true
at $\fM_w(\lambda^*_w)$ for every $w\in W_\rho$. The \emph{laws} of a frame are the
sentences true at every $\fM_w(\lambda)$, $\lambda\in\dom\fM_w$.
\end{definition}

For a projectile, $\fM(\lambda)$ solves $\ddot{\mathbf x}=-g\hat z-k|\dot{\mathbf
x}|\dot{\mathbf x}$ with parameters $(g,v_0,\theta,k)$, and ``neglect air'' deforms $k$
toward $0$.

\begin{definition}[Contexts; filter semantics]\src{T3 Defs 1.2--1.3}\label{def:physics:filter}
A \emph{context tree} has root $@$ and parent map $\pi$; each other context is
$\SUP(A)$, a supposition of a sentence $A$; or $\DEF(D,\lambda^\circ,\fG)$, a
\emph{declared deformation} with deformed parameters $D\subseteq\{1,\dots,n\}$, ideal
values $\lambda^\circ\in\R^D$, and a proper filter $\fG$ on $\R^D$ converging to
$\lambda^\circ$ (the \emph{path}); or $\AUX$, a definitional extension that shares its
parent's semantics. A $\DEF$ context uses \emph{limit semantics} if $\fG$ is principal at
$\lambda^\circ$, and \emph{germ semantics} if $\fG$ is the punctured neighbourhood filter,
possibly restricted to a curve. Fix a world $w$, write
$\tset{\varphi}_w=\{\lambda\in\dom\fM_w:\fM_w(\lambda)\models\varphi\}$ and
$\lambda[D{:=}\mu]$ for $\lambda$ with its $D$-coordinates replaced by $\mu$. Each context
$c$ gets a (possibly improper) filter $\fF_c(w)$ on $\dom\fM_w$: $\fF_@(w)$ is principal at
$\lambda^*_w$; for $c=\SUP(A)$ with parent $p$, $\fF_c(w)$ is generated by
$\fF_p(w)\cup\{\tset{A}_w\}$; for $c=\DEF(D,\lambda^\circ,\fG)$, it is generated by the sets
$\{\lambda[D{:=}\mu]:\lambda\in F,\mu\in G\}\cap\dom\fM_w$, $F\in\fF_p(w)$, $G\in\fG$.
\emph{Truth in a context} is $c\models_w\varphi$ iff $\tset{\varphi}_w\in\fF_c(w)$, and the
context-free proposition $\ist(c,\varphi)$ is true at $w$ iff $c\models_w\varphi$. A context
is \emph{proper} at $w$ if $\emptyset\notin\fF_c(w)$. A claim is \emph{supertrue} under
$\rho$ if it is true at every $w\in W_\rho$.
\end{definition}

Below a root parent, a germ context makes $\varphi$ true iff $\varphi$ holds for all
deformed parameters close enough to the ideal along the path, at the actual values of
everything else; a limit context is evaluated at the single point
$\lambda^*[D{:=}\lambda^\circ]$, and is improper if no model exists there.

\begin{proposition}[Classicality without explosion; \L{}o\'s]\src{T3 Prop 1.4}\label{prop:physics:los}
Let $\Th_w(c)=\{\varphi:c\models_w\varphi\}$. \textup{(a)} $\Th_w(c)$ is closed under
first-order consequence. \textup{(b)} It is consistent iff $c$ is proper at $w$.
\textup{(c)} $\Th_w(c)=\bigcap_U\Th(\prod_U\fM_w)$, where $U$ ranges over the
ultrafilters on $\dom\fM_w$ extending $\fF_c(w)$ and $\prod_U\fM_w$ is the ultraproduct of
the structures $\fM_w(\lambda)$.
\hfill\leanok{Contexts.cn2\_closed\_ThIn}\ \leanok{consistent\_ThIn\_iff} ((a), (b); propositional)
\end{proposition}
\begin{proof}[Proof idea]
(a), (b) by compactness and filter closure; (c) because a set lies in a filter iff it lies
in every ultrafilter extending it, with \L{}o\'s's theorem \citep{los1955quelques}. Details
in Appendix~\ref{app:physics}.
\end{proof}

If $c$ and all its $\DEF$ ancestors use limit semantics, $\fF_c(w)$ is principal or
improper, so the ultraproducts in (c) are just the standard model at the context's point.
Under germ semantics with a root parent they are nonstandard models in which the deformed
parameters are infinitely close to the ideal along the path and the undeformed ones take
\emph{exactly} their actual values: a concrete semantics for McCarthy's $\ist$, and the
sense in which Robinson vindicated Leibnizian infinitesimals
\citep{robinson1966nonstandard}. $\Th_w(c)$ may be incomplete (germs can oscillate); that
is underdetermination the reading must settle.

\begin{proposition}[Supposition is exact discharge]\src{T3 Prop 1.5}\label{prop:physics:sup}
For $c=\SUP(A)$ with parent $p$: $c\models_w\psi$ iff $p\models_w A\to\psi$. In
particular, $c$ is improper at $w$ iff $p\models_w\neg A$.
\hfill\leanok{Contexts.holdsIn\_sup\_iff} (propositional)
\end{proposition}
\begin{proof}
$X\in\fF_c(w)$ iff $X\supseteq F\cap\tset{A}$ for some $F\in\fF_p(w)$, and
$\tset{\psi}\supseteq F\cap\tset{A}$ iff $\tset{A\to\psi}\supseteq F$. Take $\psi=\bot$.
\end{proof}

\begin{proposition}[What idealized contexts may do]\src{T3 Prop 1.6, rev.}\label{prop:physics:idealized}
\textup{(a)} \emph{Contradicting the root.} In a hydrostatics frame whose laws do not
mention phase change (so the model exists at $p_{\rm atm}=0$), let $c$ deform $p_{\rm atm}$ toward $0$ (limit semantics). Then
$c\models p_{\rm atm}=0$ and $@\models p_{\rm atm}=1.0\times10^5$, both contexts are
proper, and no sentence is both true and false in one context. \textup{(b)}
\emph{Inconsistent limit, consistent germ.} A rope of mass $m_r$ is pulled by a force
$\Phi\neq0$, with $\dom=\{m_r>0\}$; deform $m_r$ toward $0$ below the root. The limit
context is improper; the germ context is proper and proves $|a|>K$ for every numeral $K$.
\textup{(c)} \emph{Norton's distinction.} Let $c_{\lim},c_{\rm germ}$ be the limit and
germ children of the root for the same $D,\lambda^\circ$, with
$\lambda^*[D{:=}\lambda^\circ]\in\dom$, and let $q(\mu)$ be the value of $Q$ at
$\lambda^*[D{:=}\mu]$. Then $c_{\rm germ}\models Q\in I$ for every open interval
$I\ni q(\lambda^\circ)$ iff $q(\mu)\to q(\lambda^\circ)$ along $\fG$ (where the deformed
models exist).
\end{proposition}

Part (c) is Norton's ``limit property versus property of the limit system''
\citep{norton2012approximation}: reasoning inside the limit model justifies a value
exactly when the query is continuous along the path; d'Alembert's paradox is the
failure case. Part (b) is the shape of idealizations that are inconsistent at the limit
but consistent along it, such as Painlevé's rigid bodies with Coulomb friction
\citep{stewart2000rigid} or a uniform isothermal gas ball in vacuum; \emph{germ semantics}
requires well-posedness along the family, not at the limit. Different filters on one
family give different contexts, so the path to the limit is part of the reading.

\paragraph{Imports are fixed by the deformation.} The import filter of an idealized
context is not a third object to learn besides rules and tolerances.

\begin{definition}[Stability]\src{T3 Def 1.7}\label{def:physics:stable}
$\varphi$ is \emph{$D$-stable} in $w$ if $\fM_w(\lambda)\models\varphi\iff\fM_w(\lambda')\models\varphi$
for all $\lambda,\lambda'\in\dom\fM_w$ differing only in coordinates in $D$. Laws are
$D$-stable, and so are \emph{undeformed data}: arithmetic formulas in the $p_i$, $i\notin D$.
\end{definition}

\begin{lemma}[Import soundness]\src{T3 Lemma 1.8, rev.: (b), (c) added}\label{lem:physics:import}
\textup{(a)} If $\varphi$ is $D$-stable in $w$ and $p\models_w\varphi$, then $c\models_w\varphi$
for every $\DEF$ child $c$ of $p$ deforming $D$, for any parent filter and path.
\textup{(b)} \emph{Uniform converse.} Fix $\fM_w$ and $D$. If for every actual point
$\lambda^*\in\dom\fM_w$ and every ideal value $\lambda^\circ\in\R^D$ the limit child $c$ of
the root satisfies $@\models\varphi\Rightarrow c\models\varphi$, then $\varphi$ is $D$-stable.
\textup{(c)} \emph{Scope.} For one fixed deformation the converse of (a) fails (e.g.\
$\varphi:=(k\le1)$ under ``$k\to0$'', limit semantics; Appendix~\ref{app:physics}).
\hfill\leanok{Contexts.import\_sound} ((a); propositional)
\end{lemma}
\begin{proof}[Proof idea]
(a) The generator built from $F=\tset{\varphi}$ lies in $\tset{\varphi}$. (b) Put the actual
point at a $\varphi$-point $\lambda$ and the ideal value at $\lambda'_D$. (c) The child always
satisfies $k=0$. Details in Appendix~\ref{app:physics}.
\end{proof}

\paragraph{Reading.} The imports that are sound \emph{uniformly} in the actual point and
the ideal value are exactly what the declared top-model family makes stable. ``Water's
vapour pressure is $2.3$\,kPa'' is a law of a frame with phase equilibrium, where the
$p_{\rm atm}\to0$ context correctly boils; in incompressible hydrostatics it is not in the
language. So learning sound import filters reduces to learning which top-model family the
problem intends: part of the reading (\S\ref{sec:physics:checker}).

% -------------------------------------------------------------------
\subsection{What coherence should constrain}\label{sec:physics:coherence}

\begin{theorem}[No truth-functional eternal reading]\src{T3 Thm 2.1, rev.}\label{thm:physics:eternalism}
Let $c$ be a $\DEF$ context, proper at $w$, with a stipulation $S_c$ such that
$c\models_w S_c$ and $S_c$ is false at $\fM_w(\lambda^*_w)$ (e.g.\ $p_{\rm atm}=0$). Let
$\tau(c,\varphi)$ translate ``$\varphi$ holds in $c$'' into a world sentence whose truth
value at $w$ is $f(S_c^w,\varphi^w)$ for some $f:\{0,1\}^2\to\{0,1\}$. Then $\tau$ does not
track in-context truth at $w$: $c\models_w\top$ and $c\not\models_w\neg S_c$, but
$\tau(c,\top)^w=\tau(c,\neg S_c)^w$.
\hfill\leanok{Contexts.not\_exists\_truthFunctional\_reading} (propositional filters)
\end{theorem}
\begin{proof}
$\top$ and $\neg S_c$ are both true at $w$, so both translations have value $f(0,1)$. And
$c\models_w S_c$ with $\Th_w(c)$ consistent (\cref{prop:physics:los}(b)).
\end{proof}

\emph{Stripping} ($f=\wedge$) makes every in-context claim false at $w$, hence incoherent
with the root; \emph{conditionalizing} ($f=\to$) makes every one true, hence vacuous; the
other two choices collapse ``in $c$'' into ``at $w$'' or its negation. Properness is
needed: for the improper rope limit of \cref{prop:physics:idealized}(b), conditionalizing
tracks in-context truth perfectly.

Consequently an \emph{eternalist} check, which deletes context indices and flags
unsatisfiable sets of recorded judgments, flags every recorded sound reductio, every sound
idealization contradicting a recorded root judgment, and every pair of sound contexts that
record different answers to one query (IPhO~2025 Theory Problems~2 and~3); what is
penalized depends on how much of the reasoning is recorded (\cref{cor:physics:eternalist},
Appendix~\ref{app:physics}).

A learner trained to keep such sets coherent loses classical logic, and not only in
the offending context. This is a formal version of the air-pressure worry of
\cref{sec:intro:question}.

\begin{proposition}[Mis-designation forces global sub-classicality]\src{T2 Prop 7.1}\label{prop:physics:misdesignation}
Let $C$ be a structural consequence relation that is $\Ac$-coherent ($A\nder_C\bot$ for
$A\in\Ac$), where some designated $A\in\Ac$ is classically inconsistent. \textup{(a)}
$C\not\supseteq\Cn_{\CPC}$; moreover, for every classical derivation of $\bot$ from $A$ in
a Hilbert or sequent-step format (no discharge metarules), some rule schema used in it
is not $C$-valid: its most general instance is not in $C$. \textup{(b)} If
$A\supseteq\{\alpha,\neg\alpha\}$, then $p,\neg p\nder_C\bot$ for every atom $p$: $C$ is
atomically paraconsistent everywhere.
\hfill\leanok{Contexts.misdesignation\_not\_Cn2\_le}\ \leanok{atomic\_paraconsistent}
\end{proposition}
\begin{proof}[Proof idea]
(a) Chaining valid instances would give $A\der_C\bot$. (b) Substitution and weakening.
Details in Appendix~\ref{app:physics}.
\end{proof}

\paragraph{Reading.} Designate ``$p_{\rm atm}=0$ plus the full background'' as coherent,
and a coherence loss teaches a sub-classical logic for \emph{all} reasoning: atomic
paraconsistency if the background states $p_{\rm atm}\neq0$ outright, the loss of some
classical schema if it only derives it. The fix has three parts. Designate only
consistent chunks \citep{brown2004chunk}. Keep the root \emph{phenomenological} (interval
claims about measurables, observations, exports), with laws living in model contexts
\citep{cartwright1983laws,wilson2006wandering}; a root holding unrestricted laws may itself
be inconsistent. And adopt \emph{meta-level eternalism}: $\ist(c,\varphi)$ is an eternal
proposition about a finitely specified context \citep{mccarthy1993notes}, by
\cref{thm:physics:eternalism} necessarily non-truth-functional in the world's facts.
Which contexts, then, must be consistent?

\subsubsection{The right constraint: realizability}

\begin{definition}[Judgment sets, bridges, realizations]\src{T3 Def 2.3, rev.}\label{def:physics:realization}
``Model'' and ``satisfiable'' refer to $\fK$; for a set $M$ of structures, $M\models\chi$
means $N\models\chi$ for all $N\in M$. A \emph{judgment set} $J$ is a finite set of triples
$(c,\Gamma,\varphi)$ over a context tree in which $\SUP$ contexts have only $\SUP$
descendants; a set $\Ccert$ of $\DEF$ contexts is \emph{certified} (claimed well-posed). A
\emph{bridge} $\beta$ has a pattern $\phi_\beta(v)$, a side condition $\sigma_\beta$ and an
export formula $\varepsilon_\beta(v)$, $v\in V$. It is a \emph{value-export bridge} for $Q$
if $\phi_\beta(v)$ is $Q=\bar v$ with standard names (distinct values denote distinct
objects in $\fK$) and $\sigma_\beta$ does not depend on $v$; it is \emph{informative} if
$\{\sigma_\beta\}\cup\{\varepsilon_\beta(v):v\in V_0\}$ has no model for some \emph{finite}
$V_0\subseteq V$ (e.g.\ interval exports over unbounded $V$). A \emph{bridge use} is
$(c,\beta,t)\in U$ with $c$ a $\DEF$ context and $(c,\emptyset,\phi_\beta(t))$,
$(\pi c,\emptyset,\sigma_\beta)$, $(\pi c,\emptyset,\varepsilon_\beta(t))\in J$. A
\emph{(frame-free) realization} assigns each context a set $M_c\subseteq\fK$ with
(R0)~$M_@\neq\emptyset$; (R1)~$M_c=M_{\pi c}\cap\tset{A}$ for $c=\SUP(A)$; (R2)~$M_c$
arbitrary, possibly empty, for $\DEF$ contexts; (R3)~every judgment holds: for
$N\in M_c$, $N\models\bigwedge\Gamma$ implies $N\models\varphi$; (R4)~every used bridge is
sound: for all $v$, $M_c\models\phi_\beta(v)$ implies $M_{\pi c}\models\sigma_\beta\to\varepsilon_\beta(v)$;
(R5)~$M_c\neq\emptyset$ for $c\in\Ccert$. The \emph{SUP-collapse} $J^*$ replaces each
$(c,\Gamma,\varphi)$ with $c$ a $\SUP$ context by $(b,\Gamma\cup\Asm(c),\varphi)$, where
$b$ is the nearest non-$\SUP$ ancestor and $\Asm(c)$ the suppositions on the path;
$\Th^*(c):=\{\bigwedge\Gamma\to\varphi:(c,\Gamma,\varphi)\in J^*\}$. The \emph{anchored set}
$\Anc$ is the least set of non-$\SUP$ contexts containing $@$ and $\Ccert$ and containing $c$
whenever $(c,\beta,t)\in U$ with $\beta$ informative and $\pi c\in\Anc$.
\end{definition}

\begin{theorem}[Realizability criterion]\src{T3 Thm 2.4, rev.: (a) frame-free; (b), (c) new}\label{thm:physics:realizability}
Suppose every used bridge is an informative value-export bridge. \textup{(a)}
\emph{Frame-free.} $J$ has a realization satisfying (R0)--(R5) iff $\Th^*(c)$ is
satisfiable for every $c\in\Anc$. \textup{(b)} \emph{Necessity in the frame semantics.}
Call $J$ \emph{frame-realizable} if for some world $w$, with the filters of
\cref{def:physics:filter}, every judgment holds ($c\models_w\bigwedge\Gamma\to\varphi$),
every $c\in\Ccert$ is proper, and every bridge use is sound at $w$ ($c\models_w\phi_\beta(v)$
implies $\pi c\models_w\sigma_\beta\to\varepsilon_\beta(v)$). If $J$ is frame-realizable,
the criterion of (a) holds, for limit or germ paths. \textup{(c)} \emph{Non-sufficiency.}
The converse of (b) fails. CE1: a certified limit child $c$ of the root deforms $D=\{2\}$
toward $0$, and $J=\{(@,\emptyset,p_1=1),(c,\emptyset,p_1=0)\}$; the criterion holds, but the
root forces $\lambda^*_1=1$, and $c$, proper because certified, sits at
$\lambda^*[2{:=}0]$, so $c\models p_1=1$. CE2: two certified siblings with the same
declaration have equal filters, so $Q=1$ in one and $Q=2$ in the other cannot both hold.
\end{theorem}
\begin{proof}[Proof idea]
By (R1), $\SUP$ contexts impose on their base exactly the constraints $\Th^*$. An
\emph{anchoring lemma} says every anchored context has $M_c\neq\emptyset$: from an empty
$M_c$ an informative bridge would export every $\varepsilon_\beta(v)$. Conversely, put one
model of $\Th^*(c)$ at each anchored context and $\emptyset$ elsewhere. Part (b) uses
properness for nonemptiness. Full proof in Appendix~\ref{app:physics}.
\end{proof}

Brute-force checks agree with (a), show that informativeness is needed, and find no
violation of (b) \status{computed} (Appendix~\ref{app:physics}). Bridge soundness here is
$\omega$-free; for judgment sets produced by the calculus of \cref{def:physics:calculus},
add $(\pi c,\emptyset,\omega_c)$ to each bridge use and read (R4) as
$M_{\pi c}\models\sigma_\beta\wedge\omega_c\to\varepsilon_\beta(v)$; the anchoring arguments
are unchanged. The locality of (a) is a \emph{consequence of (R2)}: in the frame semantics
rigid parameters and shared filters force cross-context constraints, which under limit
semantics are a pointwise adjunction (\cref{prop:physics:framerealizability}), the
structural core of discussive and preservationist logics
\citep{jaskowski1969propositional,schotch1980inference}. Germ semantics is open.

Consequences for reductio and ill-posed contexts are collected in
\cref{cor:physics:immunity} (Appendix~\ref{app:physics}): a $\SUP$ context deriving $\bot$
contributes exactly the negation of its suppositions, an idealized one deriving $\bot$
constrains frame-free realizability only if anchored, and eternalist coherence is strictly stronger than frame-free
realizability.

\paragraph{Reading.} Arguments for $P$ and $\neg P$ in contradictory contexts should not be
penalized, and \cref{thm:physics:realizability}(a) says exactly which contexts \emph{must}
be consistent: the root, and idealizations that are certified or export informatively.
Ill-posed idealizations must not export, and this is now an iff. The frame-semantic
constraints CE1, CE2 are enforced in the checker by a declaration discipline, not by
coherence (\S\ref{sec:physics:checker}).

\begin{remark}[Belief states, worlds, contextuality]\label{rem:physics:existence}
\Cref{sec:existence} gives the multi-context counterpart: coherence on designated contexts
is equivalent to a model nonempty there (\cref{thm:existence:mcs,cor:existence:contexts});
bridges must be read as rules of proof over belief states, not conditionals between worlds
(\cref{thm:existence:rules}); and locally coherent contexts need not glue into one world
(\cref{prop:existence:triangle}), a logical form of contextuality
\citep{abramsky2011sheaf}.
\end{remark}

\subsubsection{A calculus, its soundness, and blame localization}

\begin{definition}[Context calculus]\src{T3 \S2.3, rev.}\label{def:physics:calculus}
Judgments are $c\cder\varphi$; $\hat R$ is the (learned) rule set and $\mathrm{Imp}_c$ the
declared import set of $c$. \rl{Ax}~$c\cder\varphi$ for $\varphi\in\Ax(c)$, where
$\Ax(@)=\Ax_\rho$, $\Ax(\SUP(A))=\{A\}$, and $\Ax(\DEF(D,\lambda^\circ,\fG))$ is the set of
\emph{canonical stipulations}: $p_i=\lambda^\circ_i$ ($i\in D$) under limit semantics,
$p_D\in U$ for definable $U\in\fG$ under germ semantics; stipulations about undeclared
parameters are not axioms. \rl{Step}~From $c\cder\varphi_1,\dots,c\cder\varphi_k$ and an
instance $\varphi_1,\dots,\varphi_k\step\psi$ of $\hat R$, infer $c\cder\psi$.
\rl{Imp}~From $\pi c\cder\varphi$ with $\varphi\in\mathrm{Imp}_c$, infer $c\cder\varphi$.
\rl{Dis}~For $c=\SUP(A)$: from $c\cder\psi$ infer $\pi c\cder A\to\psi$. \rl{Exp}~For a
$\DEF$ context $c$: from $c\cder\phi_\beta(t)$, $\pi c\cder\sigma_\beta$ and
$\pi c\cder\omega_c$, infer $\pi c\cder\varepsilon_\beta(t)$; here the
\emph{well-posedness certificate} $\omega_c$ is a sentence such that for every $w$ and
$\lambda\in\tset{\omega_c}_w$, $\{\mu:\lambda[D{:=}\mu]\in\dom\fM_w\}\in\fG$. An instance is
\emph{sound at $w$} if it preserves truth there (for Step:
$c\models_w\bigwedge_i\varphi_i\to\psi$), and \emph{sound} if sound at every $w\in W_\rho$.
\end{definition}

Step soundness is \emph{context-relative}: time reversal is sound in a drag-free $\DEF$
context and unsound at the root.

\begin{theorem}[Soundness]\src{T3 Thm 2.5}\label{thm:physics:soundness}
If every Ax in a derivation of $c\cder\varphi$ is true at $w$ and every Step, Imp and Exp
instance in it is sound at $w$, then $c\models_w\varphi$.
\end{theorem}
\begin{proof}
Induction: Ax and Step by hypothesis and filter closure (\cref{prop:physics:los}(a));
Imp and Exp by soundness of the instance; Dis by \cref{prop:physics:sup}.
\end{proof}

\begin{corollary}[Blame localization]\src{T3 Cor 2.6, rev.}\label{cor:physics:blame}
Let $W_\rho\neq\emptyset$. Call $d$ \emph{designated} if it is the root, or a certified $\DEF$
context with a \emph{properness witness} $w_d\in W_\rho$ at which $d$ is proper. Let $\mathcal D$
derive $d\cder\bot$, and let its \emph{cone} be all instances in its derivation tree,
including those in sub-contexts whose conclusions reach $d$ by Dis or Exp. (a) At every
$w$ at which $d$ is proper (every $w\in W_\rho$ if $d=@$), the cone contains an Ax false
at $w$ or an instance unsound at $w$. (b) Items outside the cone are never implicated: an
undischarged $\SUP$ contradiction, or a pair of contradictory $\ist$-claims. (c) If
$\pi d\models_w\omega_d$ and $\pi d$ is proper at $w$, then $d$ is proper at $w$; so
parent-derived certificates are properness witnesses.
\end{corollary}
\begin{proof}[Proof idea]
(a) is the contrapositive of \cref{thm:physics:soundness}; (b) follows. (c) In each parent
filter set pick $\lambda\in\tset{\omega_d}$. \Cref{prop:physics:propagation} adds how $\bot$ in an anchored $\DEF$ context
propagates toward the root or a certified context, and where the blame then falls.
\end{proof}

The cone is a \emph{negative bag} in the sense of \cref{sec:coherence}, and by (a) a
correct target never contains it whole. So with the designated contexts of
\cref{cor:physics:blame} (the root, and certified contexts with properness witnesses) as
the \emph{only} designated contexts, an oligarchic-halving learner makes at most
$\log_2(1/\wstar)$ detections and no false alarms while the reading and the certifications
are correct (\cref{thm:coherence:halving}; false alarms are priced by
\cref{thm:coherence:robust}). What root coherence cannot see is silent unsoundness:
exports that are false but consistent with everything else (\S\ref{sec:physics:learning}).

\subsubsection{Reductio hygiene, and $100=99.9$}

The motivating notes' second worry (\S\ref{sec:physics:problem}) continues: ``and there's a
chance such an argument could equally be provided starting from the statement itself''.
Let $c$ have an explicit \emph{imported fragment} $K_c$, a finite set of premises. A
\emph{reductio of $A$ in $c$} is an $\hat R$-derivation of $\bot$ from $K_c\cup\{A\}$ in
$\SUP(A)$; accepting it yields $c\cder\neg A$. An instance is \emph{$c$-sound at $w$} if
$c\models_w\bigwedge\varphi_i\to\psi$.

\begin{theorem}[Reductio hygiene]\src{T3 Thm 1.9, rev.: (c)--(e)}\label{thm:physics:hygiene}
\textup{(a)} \emph{Soundness.} If $c\models_w K_c$ and every instance used is $c$-sound at
$w$, every accepted reductio of $A$ yields $c\models_w\neg A$. \textup{(b)}
\emph{Inconsistent imports refute everything.} If $K_c\der_{\hat R}\bot$, every sentence,
$A$ and $\neg A$ alike, has a reductio in $c$, whatever the world. \textup{(c)}
\emph{Derivation-local tests fail.} Let $K=\{A\to q,\neg A\to q,\neg q\}$. $K$ is
unsatisfiable, but the reductio of $A$ uses only the satisfiable $\{A\to q,\neg q\}$ and
that of $\neg A$ only the satisfiable $\{\neg A\to q,\neg q\}$; any test that is a function
of the derivation alone and accepts every sound reductio from a consistent premise set
accepts both. \textup{(d)} \emph{A consistency certificate suffices.} If a structure $N$
satisfies $K_c$ and every instance used ($N\models\bigwedge\varphi_i\to\psi$), every
reductio-refuted $A$ is false in $N$, so no $A$ and $\neg A$ are both refuted. The
condition on instances is automatic for logically sound $\hat R$ and needed otherwise. If
$c$ is proper at $w$, $c\models_w K_c$ and the instances are
$c$-sound at $w$, a suitable $N=\fM_w(\lambda)$ exists: a properness witness is a
certificate. \textup{(e)} \emph{Learning consequence.} A coherence loss fed with reductios
from a context whose fragment is $\hat R$-inconsistent receives ``$A$ false'' and ``$\neg A$
false'' for every $A$: jointly unsatisfiable, systematic, world-independent labels.
\hfill\leanok{Contexts.refutes\_of\_inconsistent}\ \leanok{derivation\_local\_test\_accepts\_both}\ \leanok{certificate\_not\_both} ((b)--(d))
\end{theorem}
\begin{proof}[Proof idea]
(a) $\SUP(A)$ refines $c$'s filter, so induction gives $\SUP(A)\models_w\bot$, and
\cref{prop:physics:sup} gives $\neg A$. (b) Monotonicity. (c) Truth tables. (d) Induction
inside $N$; the instance condition is needed when $c\models q$ and $K_c=\{\neg q\}$. Details
in Appendix~\ref{app:physics}.
\end{proof}

\paragraph{Reading.} A reductio in an inadequate framework can indeed be spurious, as the
notes suggest; the remedy is not to abandon reductio. Hygiene is a \emph{global}
certificate on the imported fragment, not a syntactic relevance condition on the
derivation, which (c) shows cannot work. In physics the certificate is a well-posedness
witness: an explicit or validated-numerical solution of the local model.

\begin{proposition}[The continuity asymmetry]\src{T3 Prop 1.10}\label{prop:physics:continuity}
Let premises $x_i=m_i$ hold only to tolerance, $x_i\in[m_i\pm r_i]$, with box $B$.
\textup{(a)} A \emph{direct} derivation of $y=g(x)$, $g$ continuous on $B$, errs by at most
$\sup_{x\in B}|g(x)-g(m)|$, which tends to $0$ with the tolerances by continuity of $g$.
\textup{(b)} \emph{Exact-equality} reasoning derives $\bot$ from approximately true
premises read as exact: $x=100$\,m (flat Earth) and $x=99.9$\,m (sphere), both true to
$\pm0.2$\,m, give $0=0.1$ by substitution, while as intervals they give
$x\in[99.8,100.1]\neq\emptyset$. \textup{(c)} \emph{Sound replacement.} If $A$ with the
exact laws implies $f(x)=0$, and $F$ is an inclusion-monotone interval extension of $f$
with $0\notin F(B)$, the interval premises refute $A$ soundly; the condition is
sufficient, not necessary.
\end{proposition}
\begin{proof}
(a) Definition. (b) Arithmetic. (c) Inclusion monotonicity \citep{moore1966interval} gives
$f(B)\subseteq F(B)\not\ni0$.
\end{proof}

\paragraph{Reading.} The notes ask: ``how come we usually avoid this sort of catastrophic
reasoning [$100=99.9$]? i don't think any nice `formal criterion' is going to be right?''
(\texttt{philosophy/philosophy of science/the domain of applicability of a scientific
mode(l).md}). Here is a partial answer, a sufficient condition rather than a definition of
good reasoning: ``$x=y$'' defines a closed set of data with empty interior, so
derivability of $\bot$ is a \emph{discontinuous} function of the data, while derived values
are continuous ones. This gives one precise reason behind a tentative suggestion in the
notes, that the danger of reductio in inadequate frameworks is ``maybe \dots\ some argument
in favor of trusting `constructive' reasoning more than `non-constructive' reasoning, at
least in messy domains'' (the note quoted in \S\ref{sec:physics:problem}). The fix is
interval typing with a \emph{margin} (\cref{thm:physics:chains}(c)).

% -------------------------------------------------------------------
\subsection{Export: what can justify leaving a context}\label{sec:physics:export}

Contexts let an idealized model contradict the world; exports carry its conclusions back.
No export is free: information about the idealized model alone, however complete (its jet,
its germ, finitely many values elsewhere), never justifies a finite-tolerance claim at the
actual parameter (\cref{thm:physics:nofree}). What justifies one is information of small
radius about the de-idealized family (\cref{prop:physics:radius}). A solver who may choose
contexts can then stipulate any answer unless side conditions are evaluated in the parent
and contexts are declared deformations (\cref{thm:physics:stipulation}).

Fix a one-parameter deformation (the multi-parameter case is identical) with ideal value
$0$ and actual value $\lambda^*\in(0,\Lambda]$. A family is represented by its query
function $Q:[0,\Lambda]\to\R$; a \emph{family class} is a set $\mathcal C$ of them. An
\emph{information map} $I$ gives what an export rule may use (the jet at $0$, the germ, a
restriction $Q|_S$, finitely many values). An \emph{export rule} $E$ maps
$(I(Q),\lambda^*)$ to ``abstain'' or a finite interval $[q\pm\epsilon]$, and is \emph{sound
on $\mathcal C$} if $Q(\lambda^*)\in E(I(Q),\lambda^*)$ whenever it does not abstain.

\begin{lemma}[Invisible perturbation]\src{T3 Lemma 3.2; \citealt{traub1988information}}\label{lem:physics:invisible}
If some $f$ with $f(\lambda^*)\neq0$ has $Q+Af\in\mathcal C$ and $I(Q+Af)=I(Q)$ for all
$Q\in\mathcal C$, $A\in\R$, then every export rule sound on $\mathcal C$ abstains at
$\lambda^*$ on every $Q\in\mathcal C$.
\hfill\leanok{Export.invisible\_perturbation}
\end{lemma}
\begin{proof}
An output $[q\pm\epsilon]$ on $I(Q)=I(Q+Af)$ would have to contain
$Q(\lambda^*)+Af(\lambda^*)$ for every $A$.
\end{proof}

\begin{theorem}[No free export]\src{T3 Thm 3.3, rev.: closure hypothesis}\label{thm:physics:nofree}
Every rule sound on $\mathcal C$ abstains at every $\lambda^*>0$, on every $Q\in\mathcal C$,
if (a) $\mathcal C+C^\infty[0,\Lambda]\subseteq\mathcal C$ and $I$ is the full jet at $0$;
or (b) $\mathcal C+C^\infty[0,\Lambda]\subseteq\mathcal C$ and $I$ is the germ at $0$, or
$Q|_S$ for any $S$ whose closure omits $\lambda^*$; or (c) $\mathcal C+\{\text{polynomials}\}\subseteq\mathcal C$
(e.g.\ $\mathcal C$ is the class of all restrictions of entire functions) and $I$ is the
jet up to a finite order $N$; or (d) $\mathcal C+C^\infty[0,\Lambda]\subseteq\mathcal C$ and
$I$ is any finite set of values $Q(\lambda_j)$, $\lambda_j\neq\lambda^*$. With only
$\mathcal C\supseteq C^\infty$, the conclusion still holds for $Q\in C^\infty$ in (a), (b)
and (d), and for every $Q\in\mathcal C$ in (d); in (b) it can fail outside $C^\infty$.
\hfill\leanok{Export.no\_free\_export\_germ}\ \leanok{no\_free\_export\_finite\_jet}\ \leanok{no\_free\_export\_samples}\ \leanok{no\_free\_export\_samples\_of\_poly\_subset} ((b)--(d); variants for (b), (c))
\end{theorem}
\begin{proof}[Proof idea]
\Cref{lem:physics:invisible} with the flat $e^{-1/\lambda^2}$ in (a), a $C^\infty$ bump at
$\lambda^*$ vanishing near $S$ in (b) and (d), and $\lambda^{N+1}$ in (c). Under the weaker
hypothesis, (d) perturbs the interpolating polynomial, and the germ of $\sqrt\lambda$ breaks
(b). Details in Appendix~\ref{app:physics}.
\end{proof}

For analytic families the full jet determines $Q(\lambda^*)$, but a sound use of a series
needs a remainder bound, which is information \emph{away from} $0$. This vindicates McMullin's and Laymon's demand for de-idealization information
\citep{mcmullin1985galilean,laymon1987scott}.

\begin{proposition}[Minimax export is the radius of information]\src{T3 Prop 3.4; \citealt{traub1988information}}\label{prop:physics:radius}
For an information value $i$ let $\mathcal C_i=\{Q\in\mathcal C:I(Q)=i\}$ and
$r(i)=\frac12(\sup_{\mathcal C_i}Q(\lambda^*)-\inf_{\mathcal C_i}Q(\lambda^*))$. A sound rule
with tolerance $\epsilon$ at $i$ exists iff $r(i)\le\epsilon$; the optimal one outputs the
midrange.
\hfill\leanok{Export.exists\_interval\_iff\_radius\_le}
\end{proposition}
\begin{proof}
A sound interval must contain $\{Q(\lambda^*):Q\in\mathcal C_i\}$, and the midrange interval
of radius $r(i)$ does.
\end{proof}

So an \emph{export certificate is information of small radius}: a uniform derivative
bound, a Gronwall bound, an invariance that pins $Q$, a monotone sandwich between the
ideal and a computed less-ideal value, or data near $\lambda^*$.

\begin{theorem}[Certified-neighbourhood and rate-qualified export]\src{T3 Thm 3.5}\label{thm:physics:neighbourhood}
(a) If the parent proves $\lambda^*\in U$ and a frame theorem gives
$\sup_{\lambda\in U}|Q(\fM(\lambda))-q|\le\epsilon$, the parent may soundly assert $Q\in[q\pm\epsilon]$.
(b) A germ claim $Q\sim q$ exports nothing at a fixed $\lambda^*$; with a rate
$|Q/q-1|\le C\|\lambda\|^k$ on $U$ it exports $|Q/q-1|\le C\sup_U\|\lambda\|^k$. (c) If the
problem declares a regime, the root is a germ context and asymptotic equivalence is the
correctness relation itself; for exp-log expressions in one parameter it is decidable
modulo zero-equivalence of constants \citep{gruntz1996computing,richardson1968some}.
\end{theorem}

Part (a) fixes the logical form of every side condition: a parent-certified set plus a
uniform bound over it, not a value at a guessed point. Certificates of this form include a
checkable Gronwall bound for finite-horizon exports (\cref{thm:physics:gronwall}); a rigorous range
certificate for projectiles under any drag dominated by quadratic drag
(\cref{thm:physics:projectile}), which for a steel ball certifies $R\in[9.653,10.483]$\,m
against a simulated $9.978$\,m and for a ping-pong ball is silent; and error budgets for
chains of bridges (\cref{thm:physics:chains}), under which ``$100-99.9$'' with tolerances
$0.2$ certifies not even a sign. All three are stated and proved in
Appendix~\ref{app:physics:certificates}.

\paragraph{Adversarial stipulation.} Context choice is an attack surface. Take a schema
``neglect $p_D$'': pattern $Q=\bar v$ in the child, side condition $s(\vec p)\le\eta$ for a
dimensionless smallness parameter (e.g.\ $kv_0^2/g$), conclusion $Q\in[v\pm\epsilon]$ in the
parent.

\begin{theorem}[Where side conditions live, and what contexts may be]\src{T3 Thm 3.9, rev.: (c)}\label{thm:physics:stipulation}
(a) \emph{Child-evaluated side conditions are unsound.} If Exp accepts $c\cder s\le\eta$ in
place of $\pi c\cder s\le\eta$, the export fires on every world (at $p_D=0$, $s=0$); on any
$C^\infty$-closed class this is unsound for every $\epsilon$. The ping-pong ball gets the
drag-free range $\pm5\%$, while its true range is $49\%$ short. (b) \emph{Free stipulations
are unsound even with parent evaluation.} If contexts may set undeclared parameters, a
solver stipulates $p_D=0$ and $g_c=g'$; the parent condition holds for the steel ball, and
the child's $Q=R_0(g')$ is any desired value. (c) \emph{Sufficiency.} Suppose contexts are
declared deformations with only canonical stipulations as axioms; $\sigma_\beta$ and a
well-posedness certificate $\omega_c$ are derived in the parent; and the schema is a
theorem of the frame: for all $w$, all $\lambda\in\dom\fM_w$ with
$\fM_w(\lambda)\models\sigma_\beta$ and all $v$, if
$\{\mu:\lambda[D{:=}\mu]\in\dom\fM_w,\ Q(\fM_w(\lambda[D{:=}\mu]))=v\}\in\fG$ then
$\fM_w(\lambda)\models\varepsilon_\beta(v)$. Then every accepted export is sound, for every
solver. Without $\omega_c$, (c) is false: in the rope frame of
\cref{prop:physics:idealized}(b) under limit semantics, the improper child proves $a=0$, the
schema holds vacuously, the parent proves $m_r\le\eta$, and Exp would accept $|a|\le\epsilon$
while $|a|=|\Phi|/m_r\ge|\Phi|/\eta$.
\end{theorem}
\begin{proof}[Proof idea]
(a), (b) are the constructions. For (c), the child-filter generator witnessing $Q=\bar v$,
intersected with the parent sets for $\sigma_\beta$ and $\omega_c$, puts the schema's
antecedent at every point of a parent filter set. Details in Appendix~\ref{app:physics}.
\end{proof}

\Cref{thm:physics:projectile} is such a schema (Appendix~\ref{app:physics:certificates}).
Germ contexts naturally conclude $Q\simeq v$ (standard part $v$); the analogous asymptotic
schema is sound below a principal parent given $\omega_c$
(\cref{lem:physics:asymptotic}), and under a nested germ parent it also needs uniform
convergence on a parent filter set.

\paragraph{Reading.} This is the search problem of \cref{sec:search} transplanted to
idealizations: a solver chooses contexts as a prover chooses steps. The remedy is the
analogue of cautious acceptance: contexts restricted to declared deformations with
canonical stipulations (Strevens's ``default values'' made into a syntactic discipline;
\citealt{strevens2008depth}), side conditions evaluated in the parent, and a parent-derived
well-posedness certificate supplied by the trusted catalogue (\cref{asm:physics:trusted}).

% -------------------------------------------------------------------
\subsection{Learning tolerances and validity regions}\label{sec:physics:learning}

Export certificates need tolerances and validity regions, which are natural targets for
learning. This subsection asks which signals learn them soundly against a solver who picks
the regime. Coherence among exports catches only tolerances that are too narrow
(\cref{thm:physics:coherencetol}); exchangeable calibration gives average coverage and
nothing under selection; and only a \emph{proved} regularity property turns validated
simulation of the top model into a sound certifier. A bridge schema $s$ has a
dimensionless regime parameter $\pi\in\Pi$ and an error function $e_s(\pi)$, the relative
difference between top model and idealization; the task is to learn \emph{where}
$e_s\le\tau$.

\begin{theorem}[Coherence cannot calibrate tolerances]\src{T3 Thm 4.1}\label{thm:physics:coherencetol}
Schemas $s=1,\dots,m$ export $[q_{s,i}\pm\epsilon_s]$ for true values $V_i$; a
\emph{coherence-only calibrator} chooses $\hat\epsilon$ from the discrepancies
$(q_{s,i}-q_{s',i})$ alone. If the environment class is closed under \emph{common shifts}
($(q,V)$ admissible implies $(q,V+b)$ admissible for $|b|\le\beta$: all models may share an
omitted effect), every $\hat\epsilon$ sound on the class has $\hat\epsilon_s\ge\beta$; with
unbounded shifts, $\hat\epsilon=\infty$.
\hfill\leanok{Export.coherence\_cannot\_calibrate}
\end{theorem}
\begin{proof}
$(q,V\pm\beta)$ give the same data, hence the same $\hat\epsilon$, and soundness for both signs
gives $2\beta\le2\hat\epsilon_s$.
\end{proof}

Worse, coherent tolerance vectors form an up-set, so a minimal coherent repair only ever
widens tolerances and never restricts a schema's regime, and two schemas sharing an
idealization never disagree: their agreement is common-mode, not independent evidence
\citep{levins1966strategy,odenbaugh2011buyer}. Appendix~\ref{app:physics:regions} gives
pendulum numbers for both effects \status{computed}.

The other signals are treated in Appendix~\ref{app:physics:regions}: split conformal
calibration fails under selection (\cref{thm:physics:conformal}); without regularity,
finitely many simulations certify only the queried points (\cref{thm:physics:noreg}); a
\emph{proved} Lipschitz constant or monotonicity gives a sound certifier at a matching
query cost (\cref{thm:physics:lipschitz,thm:physics:monotone}); and monotonicity can fail,
for a damped pendulum (\cref{prop:physics:laymon}).

\paragraph{Reading.} The monotonicity that \citet{laymon1987scott} uses to explicate
idealized data (better data, better predictions) is a property of a model--query pair, to
be proved, not assumed. Once such a property is proved, simulation of the top model is a
sound oracle for validity regions; world feedback proper is needed only to calibrate the
top model itself, and for olympiad grading, where the root is $c_P$, not even that.

% -------------------------------------------------------------------
\subsection{The Structured Physics Solution checker}\label{sec:physics:checker}

\begin{definition}[Structured Physics Solution]\src{T3 Def 5.1; L9 \S7.1}\label{def:physics:sps}
An \emph{SPS} is a tuple $(\rho,\mathcal T,\mathcal D,X)$: a purported reading $\rho$ (frame
class and parameter roles: given, measured, free, deformable; the query, its dimension and
allowed symbols; conventions), which, if the setter supplies a formal specification, must
be that specification or agree with it ($W_\rho$, $Q$ and its type equal, $\Ax_\rho$ its
axioms or kernel-derivable from them); a context tree $\mathcal T$; a derivation $\mathcal D$ in the
calculus of \cref{def:physics:calculus} whose Exp instances cite schemas from a catalogue
$\mathcal B$; and a final judgment $X=(@\cder Q\in[q\pm\delta])$. A declared-regime problem
uses a top-level germ context $d$ whose declaration belongs to the trusted specification,
and $X=(d\cder\chi)$.
\end{definition}

\begin{definition}[The checker]\src{T3 Def 5.2, rev.: Ax legality, schema indexing}\label{def:physics:checker}
The checker accepts iff all of the following hold. \emph{Reading:} if a formal specification is
supplied, $\rho$ is it or passes the comparison of \cref{def:physics:sps} (check C0 in
\S\ref{sec:physics:example}). \emph{Ax legality:} every Ax at the root is in $\Ax_\rho$,
every Ax in $\SUP(A)$ is $A$, every Ax in a $\DEF$ context is a canonical stipulation of its
declaration. \emph{Kernel and laws:} every Step is accepted by the kernel, and every law it
cites is in the trusted law catalogue. \emph{Schema indexing:} every Exp cites a
$\beta\in\mathcal B$ whose catalogue entry is indexed by \emph{exactly} the context's
declaration $(D,\lambda^\circ,\fG)$; $\sigma_\beta$ and the well-posedness certificate
$\omega_\beta$ listed in the catalogue entry of $\beta$ (used as $\omega_c$) are derived
\emph{in the parent}, so the solver may not choose $\omega_c$; asymptotic schemas are cited
only below a principal parent (the root, or a limit-semantics context all of whose $\DEF$
ancestors use limit semantics). \emph{Imports:} every Imp into a $\DEF$ context is a
catalogue law or an undeformed datum for its $D$ (imports into $\SUP$ children are
unrestricted). \emph{Shape:} every non-root context is $\SUP$, $\DEF$ or $\AUX$, every $\AUX$
bridge is a kernel-checked theorem, and $X$ is derived at the root (or the declared germ
root). All checks are deterministic.
\end{definition}

Ax legality and schema indexing are necessary: without them the checker accepts the attack
of \cref{thm:physics:stipulation}(b). The list of ``deformable'' parameters need not be
trusted: if a solver uses $\DEF(\{k,g\},(0,g'))$, the ``neglect $k$'' schema, indexed by
$\DEF(\{k\},0)$, cannot be cited.

\begin{assumption}[Trusted base]\src{T3 \S5, rev.}\label{asm:physics:trusted}
(TB1) The kernel and the law catalogue are sound: every accepted Step instance and every
catalogue law is valid in every $\fM_w(\lambda)$, $\lambda\in\dom\fM_w$, $w\in W_\rho$
(including top-model theorems used as certified-neighbourhood bridges). (TB2) Every
$\beta\in\mathcal B$ is a theorem of the frame for its indexed declaration, in the pointwise
form of \cref{thm:physics:stipulation}(c) or the asymptotic form of
\cref{lem:physics:asymptotic}; and each catalogue entry carries a well-posedness
certificate $\omega_\beta$ for its declaration: for every $w$ and
$\lambda\in\tset{\omega_\beta}_w$, $\{\mu:\lambda[D{:=}\mu]\in\dom\fM_w\}\in\fG$. (TB3)
Validated numerics return correct enclosures.
\end{assumption}

\begin{theorem}[Relative soundness of the SPS checker]\src{T3 Thm 5.3, rev.}\label{thm:physics:relative}
Assume (TB1)--(TB3), and that $\rho$ is a reading in the sense of
\cref{def:physics:frame}, i.e.\ $\Ax_\rho$ holds at every $w\in W_\rho$ (guaranteed when the
Reading check compares $\rho$ with a formal specification). For every solver whatsoever,
if the checker accepts $(\rho,\mathcal T,\mathcal D,X)$ with a root judgment $X$, then
$Q^{\fM_w(\lambda^*_w)}\in[q\pm\delta]$ for every $w\in W_\rho$ (for a germ root,
$d\models_w\chi$). Hence if the intended reading satisfies $W_{\rho^*}\subseteq W_\rho$ the
export is supertrue for the intended problem, and if the actual system is in
$W_{\rho^*}$ it is true of the world.
\end{theorem}
\begin{proof}[Proof idea]
\Cref{thm:physics:soundness} at each $w\in W_\rho$: Step by (TB1); root axioms because
$\rho$ is a reading; $\DEF$ axioms because canonical stipulations hold in the child filter
(the check that blocks \cref{thm:physics:stipulation}(b)); imports by
\cref{lem:physics:import}(a); exports by \cref{thm:physics:stipulation}(c) with (TB2), (TB3)
and schema indexing (by (TB2), $\omega_c$ is a genuine certificate). Nothing depends on how
$\mathcal D$ was produced. Details in Appendix~\ref{app:physics}.
\end{proof}

The residual \emph{judgment layer} is exactly (J1) \emph{reading adequacy}: $\rho$ is a
reading ($\Ax_\rho$ holds at every $w\in W_\rho$) and $W_{\rho^*}\subseteq W_\rho$; when
$\rho$ is the setter's formal specification only the second half remains; (J2)
\emph{top-model adequacy}, that the actual system is a world in $W_{\rho^*}$ (needed only
for world claims, not for olympiad grading, where $c_P$ is the root); and (J3) the
\emph{trusted base} (TB1)--(TB3). (J3) is the trusted core of any proof checker. (J1) and
(J2) cannot be removed:

\begin{theorem}[The judgment layer is ineliminable]\src{T3 Thm 5.4, rev.}\label{thm:physics:judgment}
(a) \emph{Reading.} Let $\mathcal R$ be the readings that the checker's inputs (text, figures,
SPS) do not exclude. A checker sound for every intended $\rho^*\in\mathcal R$ accepts $X$ only
if $X$ holds on $\bigcup_{\rho\in\mathcal R}W_\rho$. (b) \emph{Open world.} Extend the frame by
an unmodelled effect $\mu$, with $\fM'(\lambda,0)=\fM(\lambda)$. If, for the given slice
$\fM'(\cdot,0)$ and actual $\lambda^*$, some admissible extension has actual $\mu^*\neq0$,
and the admissible extensions are closed under $Q\mapsto Q+Af(\mu)$ for all $A\in\R$ and
$C^\infty$ $f$ with $f(0)=0$, then every finite-tolerance world claim accepted by a checker
that sees only the slice $\fM'(\cdot,0)$ and the SPS is false in some admissible extension.
\end{theorem}
\begin{proof}
(a) If $X$ fails at some $w\in W_{\rho_0}$, $\rho_0\in\mathcal R$, the checker is unsound for
$\rho^*=\rho_0$. (b) \Cref{lem:physics:invisible} in the $\mu$ direction, with $f(\mu^*)=1$.
\end{proof}
If $\mu^*=0$, or if laws bound $Q$, (b) can fail: then the open world has been closed by a
modelling assumption, i.e.\ moved into (J2).

\begin{theorem}[Supervaluational correctness; completion invariance]\src{T3 Thm 5.5, rev.}\label{thm:physics:superval}
Write $Q^w=Q^{\fM_w(\lambda^*_w)}$, with $W_\rho$ ranging over all admissible completions (free
parameters $\theta_f\in\Theta_f$, measured ones in their intervals). (a) ``$Q\in[q\pm\delta]$'' is
supertrue iff $\sup_{w\in W_\rho}|Q^w-q|\le\delta$; a supertrue answer at precision $\delta$
exists iff $\osc_{W_\rho}Q\le2\delta$, and the midrange is one. (b) A
\emph{completion-invariance certificate} is a root derivation of $Q\in[q\pm\delta]$ from
$\Ax_\rho$ in which $\theta_f$ enters only through ``$\theta_f\in\Theta_f$''. By
\cref{thm:physics:relative} the conclusion of every accepted root derivation holds at every
completion, so in the formal checker completion invariance is not a separate obligation but
the content of deriving from $\Ax_\rho$ rather than from a narrowed axiom set. Informal forms
are $\partial Q/\partial\theta_f=0$ on a connected $\Theta_f$ \emph{where $Q$ is
differentiable in $\theta_f$}, or a variation bound; in \S\ref{sec:physics:example}, $E$ is
constant in $\alpha$ only while the point stays lit for every admissible $\alpha$, and jumps
to $0$ at $s=L\tan\alpha$. (c) Supertruth over $\bigcup\mathcal R$ is the
\emph{largest} acceptance criterion sound under reading uncertainty.
\end{theorem}

\paragraph{Reading.} This answers the question quoted at the opening of this section: a
claim is true in the context at hand iff it holds on the context's filter in every
admissible completion of the reading. It is supervaluation \citep{fine1975vagueness} over
completions of an underspecified problem. Checking it decomposes into kernel steps, stable
imports, certified bridges with parent-evaluated side conditions, and a reading; only the
reading is not mechanizable. Hence our strongest recommendation: \emph{the problem setter
should supply a formal specification of the admissible completions}, after which the
checker is sound relative to the top model and the trusted base (TB1)--(TB3) alone. An
episode in the notes' introspective solution of the problem of
\S\ref{sec:physics:example} (``i suspected [the angle doesn't matter] but then i messed up a
calculation'', \texttt{introspection/introspect-solving physics olympiad
problems/introspect-solving eupho 2025-T1.md}) is a failed completion-invariance obligation
that a checker would flag at once.

\begin{proposition}[Gricean determinacy is one-sided]\src{T3 Prop 5.6}\label{prop:physics:gricean}
Call $\rho$ \emph{determinate at $\delta$} if $\osc_{W_\rho}Q\le2\delta$. If
$W_{\rho'}\subseteq W_\rho$ then $\osc_{\rho'}Q\le\osc_\rho Q$. So ``the problem is well-posed''
can refute only (some) over-wide readings, never over-narrow ones; and silent narrowing
(fixing the sun's angle, choosing a specular mirror, importing an extra law, deciding ``the
finger is horizontal'') never decreases determinacy, while it violates (J1).
\end{proposition}
\begin{proof}
A supremum over a subset is at most the supremum over the set; dually for infima.
\end{proof}

The free internal signal for readings thus pushes \emph{toward} the unsafe error: silent
narrowing is the solver's natural attack. With \cref{sec:coherence} and
\cref{thm:physics:coherencetol} this completes a table of three one-sided signals:
\begin{table}[ht]\centering\small
\caption{Three objects, three free signals, three blind spots.}\label{tab:physics:onesided}
\begin{tabular}{@{}llll@{}}
\toprule
Object & Free internal signal & Error it catches & Error it misses\\
\midrule
rules & coherence (\cref{sec:coherence}) & over-acceptance & under-acceptance\\
tolerances & coherence across exports (\cref{thm:physics:coherencetol}) & too narrow & too wide\\
readings & Gricean determinacy (\cref{prop:physics:gricean}) & (some) too wide & too narrow\\
\bottomrule
\end{tabular}
\end{table}
The opposite signal must come from outside: a formal specification by the setter (which
moves (J1) into the checked part); intended-completion data (official solutions, marking
schemes; the introspection note, on whether a finger in the part-(b) photograph is tilted:
``ok i peeked at the solution and it seems to assume that's not the case''); and
round-trip checks such as rendering the model's observable and comparing it with the
figure, which are also one-sided: passing one does not certify a reading.
\Cref{sec:philosophy} returns to the table.

% -------------------------------------------------------------------
\subsection{Worked example: EuPhO 2025 Theory 1(a), the polished chair leg}\label{sec:physics:example}

The setup is reconstructed from the introspection note of \S\ref{sec:physics:checker} and
has not been checked against the official text; the answer
$E=\frac{aI_0}{2}|\sin\frac\varphi2|/r$ is the one reached in that note, which records it as
matching the official answer. \emph{Specification:} a vertical specular cylinder of radius
$a$, lit over height $L$, stands on a flat floor; the sun is a parallel beam of irradiance
$I$ at zenith angle $\alpha$, travelling in $+x$; the floor's direct illuminance is
$I_0=I\cos\alpha$. The query is the reflected illuminance (``surplus'') $E(r,\varphi)$, with
$\varphi$ measured from the anti-sun direction $+x$. Given $a,I_0,L$; free completion
$\alpha\in[30^\circ,60^\circ]$, which the photograph does not fix; conventions: perfect
specular reflectance, geometric optics, incoherent addition, sun radius
$\delta_s=4.65$\,mrad, leg vertical; answer type: a function of $(a,I_0,r,\varphi)$ with the
dimension of $I_0$. Write $e(\vartheta)=(\cos\vartheta,\sin\vartheta)$.

\begin{proposition}[Exact reflected illuminance]\src{T3 Prop 6.1}\ \status{SymPy-verified}\label{prop:physics:legexact}
Parametrize the lit side by $\vartheta\in(\pi/2,3\pi/2)$ and the reflection height by
$h\in[0,L]$; let $\psi=2\vartheta-\pi$, $s=h\tan\alpha$ and $c=\sin\frac\psi2$. (a) The
reflected ray lands at
$P=a\,e(\vartheta)+s\,e(\psi)=(s+ac)\,e(\psi)+a\cos\frac\psi2\,e(\psi+\frac\pi2)$. (b)
$(\vartheta,s)\mapsto P$ is injective with $|\det\partial P/\partial(\vartheta,s)|=2s+ac$. (c)
$E(P)=I_0\,ac/(2s+ac)$ for $s\le L\tan\alpha$, and $E=0$ off the image and where the preimage
has $s>L\tan\alpha$; the total reflected flux $2aLI\sin\alpha$ equals the intercepted flux.
(d) At a fixed floor point $E$ does not depend on $\alpha$; only the lit domain
$s\le L\tan\alpha$ does.
\end{proposition}

Part (d) is the completion-invariance certificate, \emph{exact} in the top model.
Injectivity (the caustic of a convex cylindrical mirror is virtual) is proved in
Appendix~\ref{app:physics}; a Monte Carlo agrees to its ${\approx}0.6\%$ noise
\status{computed}. The formula is elementary catoptrics; we have not found it in the
sources checked, but the official solution may contain it, so its novelty is uncertain.

\begin{proposition}[Thin-leg bridge, rigorous]\src{T3 Prop 6.2, rev.: lit domain}\label{prop:physics:thinleg}
Let $E_0=\frac{aI_0}{2}\sigma/r$, $\sigma=|\sin\frac\varphi2|$, $\epsilon=a/r$,
$m=\frac12\arcsin\epsilon$. If $\epsilon<2/\sqrt5$, $\sigma>m$ and $r\le L\tan\alpha$ (or
$L=\infty$), then
$1-\frac m\sigma\le\frac{E}{E_0}\le\frac{1+m/\sigma}{\sqrt{1-\epsilon^2}-\epsilon/2}$.
\end{proposition}

The bound is rigorous but conservative \status{computed}; its export domain is an annulus
$a\ll r<L\tan\alpha$, an instance of intermediate asymptotics
\citep{barenblatt1996scaling}. A \emph{finite-sun bridge} (\cref{prop:physics:finitesun},
Appendix~\ref{app:physics:example}) widens it by the sun's azimuth spread, for any radiance
profile over the solar disc.

\paragraph{The checker run.} A mini-checker (\texttt{sps\_checker.py}, re-run in
\cref{sec:experiments}) encodes the specification with $L=60a$ and runs concrete versions
of the checks of \cref{def:physics:checker}: (C0) reading containment against the formal
specification; (C1) answer type: symbols $(a,I_0,r,\varphi)$ only, and dimension; (C2)
declared imports are in the specification; (C3) completion invariance: identity with the
certified $\alpha$-free form; (C4) side conditions recomputed from the root's domain, every hypothesis of
\cref{prop:physics:thinleg,prop:physics:finitesun} asserted (the script checks
$\alpha_{\max}\le\pi/2$; the stronger $\alpha_{\max}+\delta_s\le\pi/2$ holds trivially here),
and the composed error budget evaluated at the worst corner. All admissible $\alpha$ light
the disc $r\le60\tan(29.73^\circ)a=34.3a$. SymPy stands in for a proof kernel.

\begin{table}[ht]\centering\small
\caption{Mini-checker run on EuPhO~2025 Theory Problem~1(a). ``Syntactic'' marks attacks
that a purely syntactic version of checks C1, C3 and C4 accepts.}\label{tab:physics:checkerrun}
\begin{tabularx}{\textwidth}{@{}>{\raggedright\arraybackslash}p{0.47\textwidth}>{\raggedright\arraybackslash}X@{}}
\toprule
Submitted SPS & Outcome\\
\midrule
Honest: $\frac{aI_0}2\sin\frac\varphi2/r$, $r\in[20a,34a]$, $\sigma\ge0.5$, tolerance $10\%$ & \ACCEPT, certified error $\le8.88\%$ (exact point-sun error $\le2.56\%$; sun term $\le0.93\%$)\\
V0: same, tolerance $5\%$ & \REJECT: budget $0.0888>0.05$ (completeness gap: the truth is within $5\%$)\\
V1: overclaimed domain $r\ge5a$ & \REJECT: budget $0.380>0.10$\\
V2: thin-leg side condition evaluated in the child ($a:=0$), declared & \REJECT{} (\cref{thm:physics:stipulation}(a)); a silent one cannot change the budget, which C4 recomputes\\
V3: $\frac{aI_0}{2\cos\alpha}\sin\frac\varphi2/r$ (a natural slip, made and later corrected in the introspection note: $I$ for $I_0$) & \REJECT{} by C1 ($\alpha$ in answer), C3 (half-oscillation $0.42$)\\
V4: reading narrowed to $\alpha=45^\circ$ & \REJECT{} by C0, only because the specification is formal; else (J1)\\
V5: declared extra import ``$10\%$ Lambertian'' & \REJECT: import not in the specification\\
V6: $2E_{\rm thin}\cos(\alpha)^{1/1000}$ (true error ${\approx}0.5$) & \REJECT{} by C1, C3; syntactic version accepts (``$\le0.0888$'')\\
V7: $\frac{23}{20}E_{\rm thin}\cos(\alpha)^{1/50}$ (true error $0.150$) & \REJECT{} by C1, C3; syntactic version accepts\\
V8: $1.5E_{\rm thin}e^{\alpha/1000}$ & \REJECT{} by C1, C3; syntactic version accepts\\
V9: $r\ge1.05a$, $\sigma\ge1$, tolerance $0.7$ (true error $0.909$) & \REJECT{} by C4 ($\epsilon=0.952>2/\sqrt5$); syntactic version accepts (``$\le0.635$'')\\
V10: $1.09E_{\rm thin}$ (control) & \REJECT{} by C3: differs from the certified form\\
V11: undeclared $10\%$ diffuse correction $1.1E_{\rm thin}$ & \REJECT{} by C3: differs from the certified form\\
\bottomrule
\end{tabularx}
\end{table}

Plausible syntactic versions of C1, C3 and C4 accept four simple attacks (V6--V9), which
are rejected only once every check is made semantic: a small instance of the lesson of
\cref{sec:search}. The prototype is still a toy; its residual defects are listed in
Appendix~\ref{app:physics:example}.

\paragraph{What remains judgment.} The photo-to-model reading (a \emph{specular} cylinder, lit
over its full height, its unsupplied reflectance, the meaning of ``illuminance surplus'', the
$\alpha$ range) and the open world (taper, tilt, floor roughness, the seat's shadow). Each
effect could be added as a declared deformation with its own bridge, moving it from judgment
to checking; what cannot be moved is the claim that \emph{these} are all the effects the
setters intend (\cref{thm:physics:judgment}); the reading can only be tested by round
trips such as rendering the predicted floor pattern. A second, germ-root example from IPhO~2025 is in Appendix~\ref{app:physics:example}.

% -------------------------------------------------------------------
\subsection{The four logical roles of physicists' checks}\label{sec:physics:roles}

Dimensional analysis, limiting cases, symmetry and conservation are often grouped together
as ``coherence checks''. They play at least four logical roles, and only the first two are
coherence in the sense of \cref{sec:coherence}:
\begin{table}[ht]\centering\small
\caption{Logical roles of the physicist's checks.\src{L9 \S3}}\label{tab:physics:roles}
\begin{tabularx}{\textwidth}{@{}>{\raggedright\arraybackslash}p{0.2\textwidth}>{\raggedright\arraybackslash}X>{\raggedright\arraybackslash}X>{\raggedright\arraybackslash}X@{}}
\toprule
Role & Checks & What it certifies & Soundness status\\
\midrule
1. invariance of the rule set & dimensions (unit group); declared context symmetries & homogeneity, equivariance of derived equations & worst case, decidable; necessary only\\
2. in-model theorem & conservation laws & a computation inside one local model & worst case within the model; silent on exports\\
3. regression against another context & limiting cases & agreement with solved sub-models & sound only for regular limits (Norton)\\
4. export side condition & order of magnitude; ``the neglected term is small'' & a bridge at a stated tolerance & sound with explicit budgets; heuristic otherwise\\
\bottomrule
\end{tabularx}
\end{table}

\emph{Dimensions} are a type system \citep{kennedy1994dimension}, and the Buckingham $\Pi$
theorem \citep{buckingham1914physically} turns ``express in terms of'' into a type
signature; but dropping a small $\Pi$ group is a bridge (role 4), not a type fact.
\emph{Symmetry} needs uniqueness, as the notes observe (\texttt{physics/symmetry.md}): a
symmetric problem implies only a symmetric \emph{set} of solutions, and a declared
symmetry gives free negative data. \emph{Limiting cases} are valid only under Norton's
continuity (\cref{prop:physics:idealized}(c)), and \emph{order-of-magnitude} checks are
export side conditions (\cref{thm:physics:chains}). Appendix~\ref{app:physics:roles}
gives the details and the closest precedent, the Andes tutor \citep{vanlehn2005andes}.

\paragraph{What is and is not established.} Contradictory contexts are fine, and
realizability says which must be consistent. Reductio in inadequate frameworks needs a
consistency certificate, and with approximate premises a margin. Untrusted solvers are
handled for everything except the reading, whose characteristic attack, silent narrowing,
no internal signal sees. Not established (\cref{sec:open}): any theorem about \emph{learning}
readings (the deep part, J1); frame realizability under germ semantics; idealizations inside
suppositions and non-value exports; automated proofs of the regularity constants of
\cref{thm:physics:lipschitz,thm:physics:monotone}; a checker beyond one worked problem. The
mini-checker illustrates the checks on one problem and twelve variants (ten attacks, one
control and one completeness test); it is not a system, and its own soundness is argued,
not proved.
